% Curva de aprendizado (Seção 5.6.1). Dados: figuras/dados/curva.csv, gerado por
% `uv run python -m src.evaluation.export_contrastes` (validade do banco, perda
% dos logs de treino). Eixo x em escala logarítmica, porque cada ponto dobra o
% anterior; painéis separados, porque as duas séries não têm escala comum.
\begin{figure}[htb]
	\Caption{\label{fig:curva} Perda na validação e validade com 25\%, 50\% e 100\% dos pares de treino}
	\centering
	\pgfplotsset{
		painelcurva/.style={
			width=0.43\textwidth, height=5cm,
			xmin=130, xmax=860,
			xtick={167,339,672}, xticklabels={167,339,672},
			xlabel={Pares de treino},
			axis lines*=left, tick align=outside, ymajorgrids, grid style={draw=black!12},
			tick label style={font=\small}, label style={font=\small},
			yticklabel style={/pgf/number format/fixed, /pgf/number format/use comma},
		},
	}
	\begin{tikzpicture}
		\begin{axis}[xmode=log, log basis x=2, painelcurva, name=perda, ymin=0.38, ymax=0.52, ylabel={Perda na validação},
				ytick={0.40,0.44,0.48,0.52},
				yticklabel style={/pgf/number format/precision=2, /pgf/number format/fixed zerofill}]
			\addplot[corSFT, thick, mark=*] table[x=pares, y=perda, col sep=comma] {figuras/dados/curva.csv};
		\end{axis}
		\begin{axis}[xmode=log, log basis x=2, painelcurva, at={(perda.east)}, anchor=west, xshift=2.3cm,
				ymin=0, ymax=100, ylabel={Saídas válidas no XSD (\%)}]
			\addplot[corSFT, thick, mark=*] table[x=pares, y=validade, col sep=comma] {figuras/dados/curva.csv};
		\end{axis}
	\end{tikzpicture}
	\Fonte{Elaborado pelo autor}
	\Nota{X-lc25, X-lc50 e A4, com o mesmo modelo, \textit{prompt} e hiperparâmetros. Perda: a do adaptador mantido, sobre os mesmos 78 pares de validação. Validade: proporção das 138 gerações aceitas pelo esquema, 87,0\% nos três braços.}
\end{figure}
